\Subhead{6.2.1}{Априорные сведения о сигнале и шуме. Задачи интерполяции, фильтрации и экстраполяции}

Априорные сведения могут включать известную форму сигнала, случайные параметры, полосу, динамическую модель, распределение шума и корреляции. Чем сильнее реальные сведения, тем точнее возможная оценка; выдуманное априорное ограничение способно сместить результат. Аддитивная модель $y(t)=x(t)+n(t)$ удобна, но не универсальна: канал может фильтровать, искажать или нелинейно преобразовывать сообщение.

Различие трёх задач задаёт доступный интервал данных относительно оцениваемого момента $t$:
\begin{itemize}
\item фильтрация: оценить $x(t)$ по наблюдениям до $t$ включительно;
\item экстраполяция/прогноз: оценить $x(t+T)$, $T>0$, по данным до $t$;
\item интерполяция/сглаживание: оценить прошлое значение по данным, включающим более поздние относительно него наблюдения.
\end{itemize}
Сглаживание может уменьшить ошибку, поскольку использует больше информации, но требует задержки и не является причинной оценкой текущего состояния без ожидания.

При квадратичной цене ошибки оптимальная среди всех измеримых оценок — условное среднее $\hat x=\langle X|Y\rangle$. Действительно,
\begin{equation}
 \langle(X-a)^2|Y\rangle=\operatorname{Var}(X|Y)+[\langle X|Y\rangle-a]^2.
\end{equation}
Минимум достигается при равенстве второго слагаемого нулю. Максимум апостериорной плотности оптимален для иной постановки потерь; он не всегда совпадает со средним. Ограничение на линейные оценки приводит к фильтру Винера; линейно-гауссовская динамика позволяет вычислять условное среднее фильтром Калмана–Бьюси.
